\subsection{Champ \texorpdfstring{$\boldsymbol{E}$}{electrique} du a une distribution de charge quelconque}

\begin{definition}[Densite de charge] \todo{Lecture 14}
$$
\rho _{\mathrm{el}}(\boldsymbol{r})=\lim_{ \Delta V \to dV } \frac{\Delta q}{\Delta V}
$$
\end{definition}
On peut definir le champ electrique comme
$$
\boldsymbol{E}(\boldsymbol{r})=\frac{1}{4\pi\varepsilon _{0}} \iiint \frac{\rho _{\mathrm{el}}(\boldsymbol{r}')(\boldsymbol{r}-\boldsymbol{r}')}{\lVert \boldsymbol{r}-\boldsymbol{r}' \rVert^{3} }~dV'
$$
\begin{definition}[Charge de surface]
$$
\sigma _{\mathrm{el}}(\boldsymbol{r})=\lim_{ \Delta S \to dS } \frac{\Delta q}{\Delta S}.
$$
\end{definition}
On peut definir le champ electrique comme
$$
\boldsymbol{E}(\boldsymbol{r})=\frac{1}{4\pi\varepsilon_{0}}\iint _{S} \frac{\sigma _{\mathrm{el}}(\boldsymbol{r}')(\boldsymbol{r}-\boldsymbol{r}')}{\lVert \boldsymbol{r}-\boldsymbol{r}' \rVert ^{3}}~dS'
$$
\begin{example}[Utile pour plus tard; demarche pas demande]
Soit un disque de rayon $R$ uniformement charge avec $\sigma _{\mathrm{el}}$ constante. Calculons $\boldsymbol{E}$ le long de l'axe $z$, donc:
\begin{align*}
\boldsymbol{E}(0,0,z) & =\frac{1}{4\pi \varepsilon _{0}}\iint _{\mathrm{disque}} \frac{\sigma _{\mathrm{el}}(\boldsymbol{r}-\boldsymbol{r}') }{\lVert \boldsymbol{r}-\boldsymbol{r}' \rVert ^{3}}~dS',\quad \boldsymbol{r}^{\mathsf{T}}=(0,0,z) \\
\end{align*}
Par symetrie $\boldsymbol{E}(0,0,z)\parallel \boldsymbol{\hat{e}}_{z}$ :
\begin{align*}
E_{z}(0,0,z) & =\frac{1}{4\pi\varepsilon_{0}}\iint _{S} \frac{\sigma _{\mathrm{el}}(z-z')}{\lVert \begin{pmatrix}
0 & 0 & 2
\end{pmatrix}^{\mathsf{T}}-\begin{pmatrix}
x' & y' & z'
\end{pmatrix}^{\mathsf{T}} \rVert ^{3}}~dS' \\
 & =\frac{1}{4\pi \varepsilon_{0}}\iint _{S} \frac{\sigma _{\mathrm{el}}z}{\underbrace{ (x'^{2}+y'^{2}+z^{2})^{3/2} }_{ (r'^{2}-z^{2})^{3/2} }}~dS'
\end{align*}
Pour l'integration sur le disque $S$, on utilise les coordones polaires :
$$
r'^{2}=x'^{2}+y'^{2},\quad dS'=r'd\theta'dr'
$$
donc
\begin{align*}
E_{z}(0,0,r) & =\frac{1}{4\pi\varepsilon_{0}}\int _{0}^{R}\int_{0}^{2\pi} \frac{\sigma _{\mathrm{el}}z}{(r'^{2}+z^{2})^{3/2} }r'd\theta'dr' \\
 & = \frac{\sigma _{\mathrm{el}}z}{4\pi\varepsilon_{0}} 2\pi \int_{0}^{R} \frac{r'}{(r'+z^{2})^{3/2}} \, dr' \\
  & =\dots = \frac{\sigma _{\mathrm{el}}}{2\varepsilon_{0}}\left( \frac{z}{|z|}-\frac{z}{\sqrt{ z^{2}+R^{2} }} \right) 
\end{align*}

En particulier $E_{z}\to \frac{\sigma _{\mathrm{el}}}{2\varepsilon_{0}} \frac{z}{|z|}$ pour $R\to \infty$.
\end{example}

\subsection{Lignes de champ \texorpdfstring{$\boldsymbol{E}$}{electrique}}

Les lignes qui sont tangentielles a $\boldsymbol{E}$ en tout point. Ces lignes sont orientees dans la direction de $\boldsymbol{E}$. 
Les lignes de champ $\boldsymbol{\ell}(s)$ (courbe parametree par $s$) sont determinees par la condition
$$
\frac{d\boldsymbol{\ell}}{ds}\times \boldsymbol{E}(\boldsymbol{\ell}(s))=0,\quad \forall s
$$
ou de maniere plus compacte $d\boldsymbol{\ell}\times \boldsymbol{E}=0$.

\section{Loi de Gauss}

\begin{theorem}
Le flux du champ electrique a travers une surface $S$ fermee est egale a la somme des charges electriques contenues dans le volume $V$ delimite par $S$, diviseer par la permitivite du vide.
\begin{itemize}
\item Sous forme integrale : 
$$
\iint _{S} \boldsymbol{E}~d\boldsymbol{S}=\frac{1}{\varepsilon_{0}}\underbrace{ \iiint _{V} \rho _{\mathrm{el}} ~dV }_{ \text{charge totale dans }V }
$$
\item Sous forme differentielle (par le theoreme de la divergence \ref{thmdiv}]) : 
$$  
\boldsymbol{\nabla \cdot E}=\frac{\rho _{\mathrm{el}}}{\varepsilon _{0}}.
$$
\end{itemize}
\end{theorem}